% !TEX program = pdflatex
%
% cv-sablonu.tex — Burak Eltaş Executive CV şablonu
%
% Bu dosya `apply` becerisi tarafından her başvuru için kopyalanıp
% düzenlenir. Sabit kalması gereken alanlar (deneyim, eğitim, sertifikalar,
% ödüller) `profil.md` ile birebir eşleşmelidir — oradaki olgusal içerik
% burada değiştirilmez. İlana göre kişiselleştirilebilecek alanlar:
% - \profilOzet (Profesyonel Özet metni) ve \hedefUnvan (başlık altı
%   unvan) komutları — dosyanın "İLANA GÖRE KİŞİSELLEŞTİRİLECEK ALANLAR"
%   bölümünde işaretlenmiştir.
% - Sidebar'daki YETKİNLİKLER listesinin sırası (madde eklemeden/
%   çıkarmadan, `profil.md` ile aynı küme kalacak şekilde).
%
\documentclass[10pt,a4paper]{article}

\usepackage[left=0cm,right=0cm,top=0cm,bottom=0cm]{geometry}
\usepackage[utf8]{inputenc}
\usepackage[T1]{fontenc}
% Not: babel[turkish] kasıtlı olarak KULLANILMIYOR. babel[turkish] +
% hyperref + enumitem birlikte yüklendiğinde, ikinci ve sonraki özel
% (renkli işaretli) itemize ortamında "Missing \endcsname inserted"
% derleme hatası oluşuyor (doğrulanmış paket çakışması). utf8+T1 Türkçe
% karakterleri (ç, ğ, ı, İ, ö, ş, ü) doğru basmaya yeterli; kaybedilen tek
% şey Türkçe heceleme kurallarıdır ki kısa CV madde metinlerinde önemsizdir.
\usepackage{xcolor}
\usepackage{amssymb}
\usepackage{paracol}
\usepackage{enumitem}
\usepackage{changepage}
\usepackage{eso-pic}
\usepackage{hyperref}

% ---------------------------------------------------------------------
% Renkler (kaynak CV'deki lacivert / turuncu vurgu paleti)
% ---------------------------------------------------------------------
\definecolor{sidebarbg}{HTML}{1B2A4A}
\definecolor{accent}{HTML}{C9862F}
\definecolor{darktext}{HTML}{1B2A4A}
\definecolor{bodytext}{HTML}{2E2E2E}

\pagestyle{empty}
\setlength{\parindent}{0pt}

\hypersetup{colorlinks=true, urlcolor=accent, linkcolor=accent}

% Sol sütunun lacivert zemini: sabit boyutlu bir kutu/minipage İÇİNE
% sarmak yerine (bu, taşan içeriği ikinci sayfaya akıtmak yerine
% SESSİZCE KESER — doğrulanmış hata), her sayfanın arkasına eso-pic ile
% bağımsız bir dikdörtgen çiziyoruz. Böylece paracol sütunları normal
% şekilde sayfa sonunda bölünebilir. NOT: Bant her sayfada TAM sayfa
% yüksekliğinde çizilir; sidebar içeriği bir sayfada erken biterse
% (ör. 2. sayfa) bandın geri kalanı sade dekoratif kalır — kaynak
% tasarımdaki organik/değişken bant yüksekliği yerine kasıtlı bir
% sadeleştirmedir.
\AddToShipoutPictureBG{%
  \AtPageUpperLeft{\color{sidebarbg}\rule[-\paperheight]{0.32\paperwidth}{\paperheight}}%
}

% ---------------------------------------------------------------------
% ==> İLANA GÖRE KİŞİSELLEŞTİRİLECEK ALANLAR (apply becerisi tarafından) <==
% ---------------------------------------------------------------------

% Profesyonel Özet: hedef ilana göre 2-4 cümleye kısaltılabilir / anahtar
% kelime vurgusu değiştirilebilir. Olgusal rakamlar (17 yıl, 500.000+
% kullanıcı, 150+ kişi) değiştirilmez.
\newcommand{\profilOzet}{%
Çağrı merkezi operasyonları, outsource tedarikçi yönetimi ve müşteri
deneyimi alanlarında 17 yılı aşkın deneyime sahibim. Telekomünikasyon ve
fintech sektörlerinde, 500.000'den fazla aktif kullanıcıya hizmet veren
büyük ölçekli operasyonları yönettim. KPI/SLA yapılandırması, AI destekli
IVR ve dijital iletişim kanalları projelerinde uçtan uca sorumluluk
üstlendim; çok tedarikçili ekosistemlerde 150'den fazla kişilik operasyonel
kadroların yönetiminde görev aldım. Veri odaklı karar alma, süreç
iyileştirme ve kalite yönetimi alanlarında uzmanım.%
}

% Hedef unvan: ilana göre güncellenebilir (ör. "Customer Experience
% Director" ilanı için başlık buna göre uyarlanabilir), ancak
% profil.md'deki gerçek unvan geçmişiyle çelişmemelidir.
\newcommand{\hedefUnvan}{Customer Experience \& Contact Center Operations Manager}

% ---------------------------------------------------------------------
% Yardımcı komutlar
% ---------------------------------------------------------------------
\newcommand{\sidebarSection}[1]{%
  \vspace{10pt}
  {\color{accent}\bfseries\fontsize{10}{12}\selectfont #1}\\[2pt]
  {\color{white!40!sidebarbg}\rule{\linewidth}{0.4pt}}\\[4pt]
}

\newcommand{\sidebarItem}[1]{{\color{white}\small $\triangleright$~#1}\par\vspace{2pt}}

\newcommand{\mainSection}[1]{%
  \vspace{8pt}
  {\color{darktext}\bfseries\fontsize{12}{14}\selectfont #1}\\
  {\color{accent}\rule{\linewidth}{1pt}}\\[4pt]
}

\newcommand{\job}[4]{%
  % #1 = unvan, #2 = şirket, #3 = tarih aralığı, #4 = madde listesi (itemize içi)
  {\color{darktext}\bfseries #1}\\
  {\color{accent}\itshape #2 \hfill #3}\\[2pt]
  \begin{itemize}[leftmargin=12pt,itemsep=1pt,topsep=2pt,label=\textcolor{accent}{$\blacktriangleright$}]
  #4
  \end{itemize}
  \vspace{4pt}
}

\begin{document}

\columnratio{0.32}
\begin{paracol}{2}

% =======================================================================
% SOL SÜTUN — SIDEBAR
% =======================================================================
\begin{leftcolumn}
\color{white}
\begin{adjustwidth}{14pt}{10pt}
\vspace*{24pt}

\sidebarSection{İLETİŞİM}
\sidebarItem{İstanbul, Türkiye}
\sidebarItem{+90 532 555 32 80}
\sidebarItem{eltas.burak@gmail.com}
\sidebarItem{\href{https://linkedin.com/in/burakeltas}{linkedin.com/in/burakeltas}}

\sidebarSection{YETKİNLİKLER}
% ==> apply becerisi bu listeyi ilanın anahtar kelimelerine göre
%     yeniden sıralayabilir; madde eklemez/çıkarmaz (profil.md ile eşleşmeli).
\sidebarItem{Customer Experience Management}
\sidebarItem{Contact Center Operations}
\sidebarItem{Outsource Vendor Management}
\sidebarItem{CRM Strategy \& Execution}
\sidebarItem{SLA \& KPI Management}
\sidebarItem{Process Improvement}
\sidebarItem{Operational Excellence}
\sidebarItem{Quality Assurance \& Audit}
\sidebarItem{Team Leadership (150+)}
\sidebarItem{Performance Analytics}
\sidebarItem{AI-Supported IVR / Digital Channels}
\sidebarItem{Capacity Planning}
\sidebarItem{Vendor Budget Management}
\sidebarItem{Eğitim Tasarımı (Train the Trainer)}

\sidebarSection{DİLLER}
\sidebarItem{Türkçe --- Ana Dil}
\sidebarItem{İngilizce --- Temel}

\sidebarSection{EĞİTİM}
{\color{white}\small\bfseries Lisans --- İşletme}\\
{\color{white!85!black}\small Anadolu Üniversitesi}\\
{\color{white!70!black}\small 2013 -- 2016}\\[6pt]
{\color{white}\small\bfseries Ön Lisans --- Bilgisayar Teknolojisi ve Programlama}\\
{\color{white!85!black}\small Çanakkale Onsekiz Mart Üniversitesi}\\
{\color{white!70!black}\small 2006 -- 2008}

\sidebarSection{ÖDÜLLER}
{\color{white}\small\bfseries Süper İş Ödülü}\\
{\color{white!85!black}\small Turkcell Global Bilgi \textbar\ 2013}\\
{\color{white!70!black}\small Üstün başarı performansı ile sertifikalı kurumsal ödül.}

\sidebarSection{SERTİFİKALAR}
\sidebarItem{Eğitimcinin Eğitimi --- BNB Training}
\sidebarItem{Temel Liderlik Becerileri --- Turkcell Akademi}
\sidebarItem{Temel Proje Yönetimi --- Turkcell Akademi}
\sidebarItem{Kök Neden Analizi --- Turkcell Akademi}
\sidebarItem{İş Süreçleri Yönetimi --- İst. İşletme Enst.}
\sidebarItem{ISO 9001 Kalite Yönetim Sistemi}
\sidebarItem{Müşteri İlişkileri Yönetimi}

\sidebarSection{DİĞER}
\sidebarItem{Kurumsal Müşteri Yönetimi}
\sidebarItem{İkna ve Müzakere Becerileri}
\sidebarItem{MS Office İleri Düzey (Excel, PPT, Word)}

\sidebarSection{GÖNÜLLÜ FAALİYET}
{\color{white}\small\bfseries Genç TEMA Üyesi}\\
{\color{white!85!black}\small Genç TEMA Topluluğu \textbar\ 2014 -- 2026}\\
{\color{white!70!black}\small TEMA Vakfı projelerinde aktif görev.}

\vspace{20pt}
\end{adjustwidth}
\end{leftcolumn}

% =======================================================================
% SAĞ SÜTUN — ANA İÇERİK
% =======================================================================
\begin{rightcolumn}
\begin{adjustwidth}{18pt}{12pt}
\vspace*{24pt}

{\color{darktext}\bfseries\fontsize{22}{24}\selectfont BURAK ELTAŞ}\\[2pt]
{\color{accent}\fontsize{12}{14}\selectfont \hedefUnvan}

\mainSection{PROFESYONEL ÖZET}
{\small\color{bodytext}\profilOzet}

\mainSection{PROFESYONEL DENEYİM}

\job{Dış Kaynak Yöneticisi}{Paribu Teknoloji A.Ş.}{2025 -- 2026}{%
\item 500.000'den fazla aktif fintech kullanıcısına hizmet veren outsource çağrı merkezi ekosisteminin uçtan uca koordinasyonunu yönettim.
\item Çok tedarikçili yapıda ortak bir performans skorlama ve raporlama sistemi kurarak tedarikçi değerlendirme süreçlerini standardize ettim.
\item AI destekli IVR ve dijital iletişim kanalı entegrasyon projelerini yöneterek self-servis kullanımını artırmaya yönelik dönüşümü hayata geçirdim.
\item Tedarikçi maliyet yapısını ve sözleşme koşullarını yeniden yapılandırarak bütçe disiplini ve maliyet optimizasyonu sağladım.
}

\job{Müşteri Hizmetleri Yöneticisi}{Paribu Teknoloji A.Ş.}{2020 -- 2025}{%
\item 5 outsource tedarikçi ve 150'den fazla operasyonel personelden oluşan müşteri hizmetleri ekosistemini yönettim.
\item KPI setlerini sıfırdan tasarladım; aylık performans denetimlerini ve SLA takip mekanizmalarını kurumsallaştırdım.
\item Outsource operasyon bütçesini ve aylık fatura mutabakat süreçlerini yöneterek finansal kontrol düzenini kurdum.
\item Kapasite planlama modeli oluşturarak yoğun dönemlerde hizmet sürekliliğini ve müşteri memnuniyetini güvence altına aldım.
\item Süreç iyileştirme ve kalite izleme projeleri yürüterek operasyonel verimlilik ve hizmet kalitesinde sürdürülebilir gelişim sağladım.
}

\job{Çağrı Merkezi Takım Lideri}{Turkcell Global Bilgi}{2017 -- 2019}{%
\item Home-agent ve sahada görev yapan 40'tan fazla kişilik ekibin işe alım, eğitim ve performans yönetimini yürüttüm.
\item Koçluk ve geri bildirim programları tasarlayarak çalışan bağlılığını ve ekip devamlılığını güçlendirdim.
\item Müşteri deneyimi sapma analizleri yürüterek ilk aramada çözüm (FCR) süreçlerinin iyileştirilmesine liderlik ettim.
}

\job{Kalite Geliştirme Uzmanı}{Turkcell Global Bilgi}{2013 -- 2017}{%
\item Finans projeleri kapsamında aylık 8.000'den fazla çağrının kalite analiz sürecini yönettim.
\item Müşteri temsilcilerine yönelik eğitim içerikleri tasarladım ve uyguladım.
\item Gizli müşteri (mystery shopping) ve kalibrasyon programları kurarak segment bazlı kalite ölçümünü standardize ettim.
\item Kök neden analizine dayalı şikâyet yönetimi yaklaşımıyla tekrarlayan şikâyetlerin önlenmesine yönelik aksiyon planlarını yürüttüm.
}

\job{Eğitim ve Kalite Sorumlusu}{Koç Sistem / Callus Çağrı Merkezi}{2010 -- 2012}{%
\item Aygaz, Vodafone, Yapı Kredi, Ford, Opet ve Dell projelerinin 7/24 inbound ve outbound operasyonlarında kalite standartlarını yönettim.
\item İşe alım mülakatlarını kalite kriterleri bazında yapılandırarak yeni temsilcilerin adaptasyon süreçlerini hızlandırdım.
\item Düşük performanslı projelerde acil aksiyon planları devreye alarak performans toparlanma süreçlerini yönettim.
}

\job{Çağrı Merkezi Müşteri Temsilcisi}{Koç Sistem / Callus Çağrı Merkezi}{2010}{%
\item Vodafone inbound/outbound operasyonlarında görev aldım; gösterdiğim performansla 9 ay içinde Eğitim ve Kalite Sorumlusu pozisyonuna terfi ettim.
}

\job{Kurumsal Satış Sorumlusu}{Avea Telekomünikasyon A.Ş.}{2008 -- 2010}{%
\item Kurumsal müşteri portföyü geliştirme ve bölge bazlı yeni müşteri kazanımı çalışmalarını yürüttüm.
\item Türk Telekom ortak satış projelerinde görev aldım; bayi denetim ve eğitim programlarıyla satış kanalı süreçlerini standardize ettim.
}

\mainSection{ANAHTAR BAŞARILAR}
\small
\begin{itemize}[leftmargin=12pt,itemsep=2pt,topsep=2pt,label=\textcolor{accent}{$\blacktriangleright$}]
\item 500.000+ aktif kullanıcıya hizmet veren outsource operasyon ekosisteminin uçtan uca yönetimi
\item 5 tedarikçi ve 150+ kişilik çok-vendorlu kadronun liderliği
\item AI destekli IVR ve dijital kanal dönüşüm projelerinin hayata geçirilmesi
\item KPI/SLA yapısının sıfırdan kurulması ve kurumsallaştırılması
\item Süper İş Ödülü --- Turkcell Global Bilgi (2013)
\end{itemize}

\vspace{16pt}
{\small\color{black!60} Burak Eltaş \textperiodcentered\ Executive CV}

\end{adjustwidth}
\end{rightcolumn}

\end{paracol}
\end{document}
